% =====================================================================
%  CARNET D'ENTRAÎNEMENT — Fiche 12 (Chimie organique)
%  THÈME 1 — Hydrocarbures, formules brutes et squelette carboné
%  NIVEAU MOYEN — CORRIGÉ
% =====================================================================
\documentclass[11pt,a4paper]{article}
\usepackage[utf8]{inputenc}
\usepackage[T1]{fontenc}
\usepackage{lmodern}
\usepackage[french]{babel}
\usepackage{amsmath,amssymb,mathtools}
\providecommand{\degree}{^\circ}
\usepackage{icomma}
\usepackage{siunitx}
\sisetup{output-decimal-marker={,},per-mode=symbol}
\usepackage[version=4]{mhchem}
\usepackage{chemfig}
\setchemfig{atom sep=2em,bond length=0.9em,cram width=2.5pt}
\usepackage{xcolor}
\usepackage{array}
\usepackage{booktabs}
\usepackage{tabularx}
\usepackage[most]{tcolorbox}
\usepackage{enumitem}
\usepackage{tikz}
\usetikzlibrary{arrows.meta,positioning,calc}
\usepackage{fancyhdr}
\usepackage{titlesec}
\usepackage[hidelinks]{hyperref}
\usepackage[a4paper,margin=2.1cm,headheight=26pt,headsep=9pt]{geometry}

\definecolor{primary}{RGB}{150,98,162}
\definecolor{primarydark}{RGB}{92,52,110}
\definecolor{excol}{RGB}{0,135,90}
\definecolor{attncol}{RGB}{200,40,55}
\definecolor{methcol}{RGB}{90,90,100}

\newcommand{\runtitle}{Thème 1 — Moyen — Corrigé}
\pagestyle{fancy}\fancyhf{}
\fancyhead[L]{\small\textcolor{primarydark}{\textbf{Fiche 12} — Chimie organique}}
\fancyhead[R]{\small\textcolor{primarydark}{\runtitle}}
\fancyfoot[C]{\small\thepage}
\renewcommand{\headrulewidth}{0.8pt}
\renewcommand{\headrule}{\hbox to\headwidth{\color{primary}\leaders\hrule height \headrulewidth\hfill}}
\setlength{\parindent}{0pt}\setlength{\parskip}{3pt}

\tcbset{boxbase/.style={enhanced,breakable,boxrule=0.9pt,arc=2.5pt,
   left=3mm,right=3mm,top=2mm,bottom=2mm,fonttitle=\bfseries,coltitle=white,
   toptitle=0.5mm,bottomtitle=0.5mm}}
\newtcolorbox[auto counter]{corr}[1]{boxbase,colback=excol!5,colframe=excol,
   title={\textsc{Corrigé}~\thetcbcounter\ \textmd{--- \itshape #1}}}
\newcommand{\rep}[1]{\textcolor{primarydark}{\textbf{#1}}}

\newcommand{\banner}[2]{%
\begin{tcolorbox}[enhanced,colback=excol,colframe=excol,arc=5pt,boxrule=0pt,
   left=5mm,right=5mm,top=2.5mm,bottom=2.5mm,coltext=white]
{\large\bfseries #1}\\[1pt]{\small #2}
\end{tcolorbox}\vspace{2pt}}

\begin{document}

\banner{Thème 1 — Hydrocarbures, formules brutes et squelette}{Niveau \textbf{Moyen} — Corrigé détaillé \quad$\vert$\quad Fiche 12, Chimie organique}

% ---------------------------------------------------------------------
\begin{corr}{du degré d'insaturation aux structures possibles}
\textbf{a)} $I=\dfrac{2\cdot5+2-8}{2}=\dfrac{4}{2}=\rep{2}$.\\
\textbf{b)} $I=2$ signifie \rep{2 insaturations} au total, c'est-à-dire toute combinaison de \emph{liaisons $\pi$} et de \emph{cycles} sommant à 2 (ex.\ une triple liaison, ou deux doubles, ou une double + un cycle, ou deux cycles).\\
\textbf{c)} Trois structures compatibles avec $\ce{C5H8}$ :
\begin{itemize}[leftmargin=1.6em,itemsep=1pt]
  \item \rep{pent-1-yne} $\ce{HC#C-CH2-CH2-CH3}$ (une triple liaison $=2$ insaturations) ;
  \item \rep{penta-1,3-diène} $\ce{CH2=CH-CH=CH-CH3}$ (deux doubles liaisons) ;
  \item \rep{cyclopentène} (un cycle + une double liaison).
\end{itemize}
\end{corr}

% ---------------------------------------------------------------------
\begin{corr}{trier les carbones par nombre d'hydrogènes}
\textbf{a)} On compte 6 carbones (chaîne « pent » + 1 méthyle). En complétant chaque C à 4 liaisons, on trouve 12 H : \rep{\ce{C6H12}}. Vérification : $I=\dfrac{2\cdot6+2-12}{2}=1$, cohérent avec l'unique double liaison.\\
\textbf{b)} En parcourant $\ce{CH2}{=}\ce{CH}{-}\ce{CH2}{-}\ce{CH}({-}\ce{CH3}){-}\ce{CH3}$ :
\begin{itemize}[leftmargin=1.6em,itemsep=1pt]
  \item \rep{3 H} : les 2 groupes $\ce{CH3}$ (bout de chaîne + méthyle ramifié) $\Rightarrow$ \rep{2 carbones} ;
  \item \rep{2 H} : le $\ce{=CH2}$ terminal \emph{et} le $\ce{CH2}$ interne $\Rightarrow$ \rep{2 carbones} ;
  \item \rep{1 H} : le $\ce{=CH-}$ de la double liaison \emph{et} le $\ce{CH}$ ramifié $\Rightarrow$ \rep{2 carbones} ;
  \item \rep{0 H} : \rep{aucun}.
\end{itemize}
Remarque : le carbone $\ce{=CH-}$ ne porte qu'\emph{un} H car sa double liaison compte pour 2 liaisons ($2+1_{\text{(vers C3)}}+1_{\text{(H)}}=4$).
\end{corr}

% ---------------------------------------------------------------------
\begin{corr}{la stéréoisomérie cis / trans}
\textbf{a)} Pour qu'une double liaison donne des isomères \emph{cis}/\emph{trans}, il faut que \emph{chaque} carbone de la double liaison porte \textbf{deux substituants différents}. Dans le but-2-ène, chaque C de $\ce{C=C}$ porte un $\ce{CH3}$ \emph{et} un H (deux groupes différents) $\Rightarrow$ deux formes. Dans le but-1-ène, le carbone terminal $\ce{=CH2}$ porte \textbf{deux H identiques} $\Rightarrow$ pas de distinction géométrique. \rep{C'est la présence de deux substituants différents de part et d'autre de C=C qui crée le cis/trans.}\\
\textbf{b)}
\begin{center}
\begin{tabular}{ccc}
\textbf{cis (Z)} & \qquad\qquad & \textbf{trans (E)} \\[6pt]
\chemfig{C(-[:150]CH_3)(-[:210]H)=[:0]C(-[:30]CH_3)(-[:-30]H)} & &
\chemfig{C(-[:150]CH_3)(-[:210]H)=[:0]C(-[:30]H)(-[:-30]CH_3)} \\[6pt]
les deux $\ce{CH3}$ du \rep{même côté} & & les deux $\ce{CH3}$ de \rep{part et d'autre} \\
\end{tabular}
\end{center}
\textbf{c)} Oui, même formule brute $\ce{C4H8}$, mais \rep{ce n'est pas la même molécule} : la liaison $\pi$ \textbf{bloque la rotation} autour de $\ce{C=C}$, donc les deux arrangements sont figés et ont des propriétés physiques différentes (points d'ébullition distincts).
\end{corr}

% ---------------------------------------------------------------------
\begin{corr}{les trois isomères de \ce{C5H12}}
\textbf{a)} et \textbf{b)} :
\begin{center}\renewcommand{\arraystretch}{1.4}
\begin{tabular}{c c l}
\toprule
Nom & Structure & Décompte des carbones \\
\midrule
\textbf{n-pentane} & \scalebox{0.7}{\chemfig{CH_3-[:30]CH_2-[:-30]CH_2-[:30]CH_2-[:-30]CH_3}} & \rep{2 primaires, 3 secondaires} \\
\textbf{2-méthylbutane} & \scalebox{0.7}{\chemfig{CH_3-[:30]CH(-[:-90]CH_3)-[:-30]CH_2-[:30]CH_3}} & \rep{3 primaires, 1 secondaire, 1 tertiaire} \\
\textbf{2,2-diméthylpropane} & \scalebox{0.7}{\chemfig{C(-[:30]CH_3)(-[:-30]CH_3)(-[:90]CH_3)-[:-90]CH_3}} & \rep{4 primaires, 1 quaternaire} \\
\bottomrule
\end{tabular}
\end{center}
\end{corr}

% ---------------------------------------------------------------------
\begin{corr}{reconstruire une molécule à partir de contraintes}
\textbf{a)} Deux carbones tertiaires liés l'un à l'autre : chacun porte donc, en plus de l'autre carbone tertiaire, \textbf{deux} autres carbones. Comme il n'y a aucun carbone secondaire, ces voisins ne peuvent être que des $\ce{CH3}$ (carbones primaires). On obtient $(\ce{CH3})_2\ce{CH-CH}(\ce{CH3})_2$ :
\begin{center}\chemfig{CH_3-[:-30]CH(-[:-90]CH_3)-[:30]CH(-[:90]CH_3)-[:-30]CH_3}\end{center}
\textbf{b)} C'est le \rep{2,3-diméthylbutane}. Décompte : \rep{2 carbones tertiaires} ($\ce{CH}$, 1 H chacun) $+$ \rep{4 carbones primaires} ($\ce{CH3}$, 3 H chacun). Total : $2+4=6$ C et $2\times1+4\times3=14$ H, soit bien $\rep{\ce{C6H14}}$. $\checkmark$
\end{corr}

\end{document}
